\chapter{Multiplication~: premiers pas}\label{ch:g2:mult}

Quatre bo\^ites de trois g\^ateaux~: plut\^ot que
$3 + 3 + 3 + 3$, on \'ecrit $4 \times 3$. La multiplication est le
raccourci pour additionner plusieurs fois le m\^eme nombre --- et sa
meilleure image est un joli tableau de lignes et de colonnes.

\section{Ajouter le m\^eme nombre}

\begin{definition}[Multiplication]\label{def:g2:mult:def}
$4 \times 3$ signifie \guillemotleft{} quatre fois trois
\guillemotright{}~:
\[
4 \times 3 = 3 + 3 + 3 + 3 = 12 .
\]
Le signe est $\times$, lu \guillemotleft{} fois \guillemotright{}.
Le r\'esultat est le \emph{produit} (la multiplication grandit dans
\cref{ch:g3:mult}).
\end{definition}

\begin{omfigure}
\begin{tikzpicture}[scale=0.5]
  \foreach \r in {0,1,2,3}
    \foreach \c in {0,1,2}
      \fill[omDef] (\c,-\r) circle (0.28);
  \node[right, font=\small] at (3.0,-1.5)
    {$4$ lignes de $3$ points~: $4 \times 3 = 12$};
\end{tikzpicture}

{\small Un \emph{tableau} (ou \emph{disposition en lignes})~: $4$
lignes de $3$. Compter par lignes donne $4 \times 3$~; les m\^emes
points par colonnes donnent $3 \times 4$ --- donc l'ordre n'a pas
d'importance.}
\end{omfigure}

\begin{example}[L'ordre n'a pas d'importance]\label{ex:g2:mult:order}
$4 \times 3$ et $3 \times 4$ comptent les m\^emes $12$ points, l'un
par lignes, l'autre par colonnes. Tourner le tableau sur le c\^ot\'e
le prouve. Donc apprendre un produit donne son jumeau gratuitement.
\end{example}

\section{Compter par bonds et les tables faciles}

\begin{method}[Construire une table en comptant par bonds]\label{met:g2:mult:skip}
Pour trouver $5 \times 4$, compter de $4$ en $4$ cinq fois~:
$4, 8, 12, 16, 20$. Compter par bonds sur la droite num\'erique
\emph{est} la table.
\begin{itemize}
  \item Table de $2$~: $2, 4, 6, 8, 10, \dots$ (les doubles)~;
  \item table de $5$~: $5, 10, 15, 20, 25, \dots$ (se termine par
        $5$ ou $0$)~;
  \item table de $10$~: $10, 20, 30, \dots$ (ajouter un z\'ero).
\end{itemize}
\end{method}

\begin{omfigure}
\begin{tikzpicture}[scale=0.6]
  \draw[-Stealth] (-0.4,0) -- (11.4,0);
  \foreach \i in {0,...,11} \draw (\i,-0.1) -- (\i,0.1);
  \foreach \i in {0,2,4,6,8,10}
    \node[below, font=\footnotesize] at (\i,-0.12) {\i};
  \foreach \i in {0,1,2,3,4}
    \draw[-Stealth, omDef, thick] (2*\i,0.28) .. controls
      (2*\i+1,0.85) .. (2*\i+2,0.28);
  \node[omDef, font=\small] at (5.5,1.1) {sauts de $2$};
\end{tikzpicture}

{\small La table de $2$ comme sauts sur la droite~: $2, 4, 6, 8,
10$. Cinq sauts de $2$ atteignent $10$, donc $5 \times 2 = 10$.}
\end{omfigure}

\begin{example}[Doubles et moiti\'es]\label{ex:g2:mult:double}
Multiplier par $2$, c'est doubler~: $2 \times 7 = 14$. L'inverse,
prendre la moiti\'e, d\'efait le double~: la moiti\'e de $14$ est
$7$. Les doubles connus par c{\oe}ur rendent toute la table de $2$
instantan\'ee --- et pr\'eparent la division
(\cref{ch:g3:division}).
\end{example}

\section{Multiplier dans les probl\`emes}

\begin{example}\label{ex:g2:mult:problem}
\guillemotleft{} Une bo\^ite contient $5$ {\oe}ufs. Combien
d'{\oe}ufs dans $4$ bo\^ites~? \guillemotright{} Quatre groupes de
cinq~: $4 \times 5 = 20$ (compter par bonds~: $5, 10, 15, 20$). Il y
a $20$ {\oe}ufs.
\end{example}

\section{Exercices}

\begin{exercise}[$\star$]\label{exo:g2:mult:1}
\'Ecris sous forme de multiplication, puis calcule~:
$2 + 2 + 2$~; \quad $5 + 5 + 5 + 5$~; \quad $10 + 10 + 10$.
\end{exercise}

\begin{exercise}[$\star$]\label{exo:g2:mult:2}
Dessine un tableau de $3$ lignes de $5$ points. Quelle multiplication
montre-t-il~? Et tourn\'e sur le c\^ot\'e~?
\end{exercise}

\begin{exercise}[$\star$]\label{exo:g2:mult:3}
Compte par bonds pour remplir la table de $2$ jusqu'\`a
$2 \times 6$~; la table de $5$ jusqu'\`a $5 \times 6$.
\end{exercise}

\begin{exercise}[$\star$]\label{exo:g2:mult:4}
Calcule~: $2 \times 7$~; \quad $5 \times 3$~; \quad
$10 \times 4$~; \quad $2 \times 9$~; \quad $5 \times 6$.
\end{exercise}

\begin{exercise}[$\star$]\label{exo:g2:mult:5}
Calcule les doubles~: le double de $6$~; le double de $8$~; le
double de $10$. Puis les moiti\'es~: la moiti\'e de $12$~; la
moiti\'e de $20$.
\end{exercise}

\begin{exercise}[$\star$]\label{exo:g2:mult:6}
Calcule la table de $3$ en comptant par bonds~: $3, 6, 9, ?, ?, ?$.
Que vaut $3 \times 5$~?
\end{exercise}

\begin{exercise}[$\star$]\label{exo:g2:mult:7}
\guillemotleft{} Il y a $4$ tables avec $5$ chaises chacune. Combien
de chaises~? \guillemotright{} \'Ecris la multiplication et la
r\'eponse.
\end{exercise}

\begin{exercise}[$\star$]\label{exo:g2:mult:8}
Une araign\'ee a $8$ pattes. Combien de pattes ont $3$ araign\'ees~?
(Compte de $8$ en $8$~: $8, 16, \dots$)
\end{exercise}

\begin{exercise}[$\star$]\label{exo:g2:mult:9}
Deux fa\c{c}ons d'arranger $12$ points en lignes \'egales~: dessine
$2$ lignes de $6$, et $3$ lignes de $4$. \'Ecris la multiplication
pour chacune.
\end{exercise}

\begin{exercise}[$\star\star$]\label{exo:g2:mult:10}
Les g\^ateaux arrivent par bo\^ites de $5$. Mia a besoin de $18$
g\^ateaux pour une f\^ete. Combien de bo\^ites doit-elle acheter~?
(Compte de $5$ en $5$ jusqu'\`a d\'epasser $18$.)
\end{exercise}
